\chapter{Unresolved mixing, stochastic forcing and diagnostic interpretation}
\label{mixing-theory}

The resolved equations do not describe every turbulent eddy, surface exchange
or particle-scattering process inside a cell. Closures supply those missing
relationships. They are mathematical models with their own state, scales and
admissible inputs, not interchangeable corrections to a finished solution.
This chapter follows the implemented Nimbus closures and their connections to
transport and diagnostics. The earlier grid and microphysics chapters supply
the carrier and moment conventions; the architecture guide explains the
corresponding state and execution owners.

\section{Eddy viscosity connects resolved strain to unresolved transport}

For physical Cartesian velocity components, define the symmetric strain
$S_{ij}=\tfrac12(\partial_j u_i+\partial_i u_j)$ and its magnitude
$|S|=(2S_{ij}S_{ij})^{1/2}$. Strain has units s$^{-1}$. An eddy viscosity
models an unresolved momentum exchange through a stress proportional to
resolved deformation. A viscous-traction form is $\tau_{ij}=2\rho_d\nu_tS_{ij}$,
with acceleration from its divergence divided by the appropriate carrier.
This motivates the closure; its isotropic/compressional and directional
treatment must still follow the actual discrete tensor owner.

Static Smagorinsky uses
\[
 \nu_t=(C_s\Delta)^2|S|,\qquad
 K_\theta=\nu_t/\mathrm{Pr}_t,\qquad
 K_q=\nu_t/\mathrm{Sc}_t.
\]
Here $C_s$ is dimensionless, $\Delta$ is a length, and viscosity/diffusivities
are in m$^2$ s$^{-1}$. Nimbus additionally applies configured coefficient
bounds and the selected stability correction. Increasing $C_s$ or the filter
length increases modeled mixing; it is not equivalent to changing a plotting
resolution. The retained \texttt{dx} filter option actually uses
$\sqrt{\Delta x\Delta y}$, while \texttt{cubic\_root} uses
$(\Delta x\Delta y\Delta z)^{1/3}$. A user filter is a prescribed physical
length. Mapped filters use the horizontal footprint and physical Jacobian
thickness, not a sloping face-normal CFL distance.

Directional modes select which gradients and exchanges participate. In the
flat derivative limit, the strain measures have the forms
\[
\begin{aligned}
 |S|_h^2&=2(S_{xx}^2+S_{yy}^2)+4S_{xy}^2
                    +(\partial_xw)^2+(\partial_yw)^2,\\
 |S|_v^2&=(\partial_zu)^2+(\partial_zv)^2+2S_{zz}^2.
\end{aligned}
\]
The horizontal-only mode therefore retains horizontal gradients of vertical
velocity. It is not simply a two-component two-dimensional atmosphere.
The C-grid implementation averages appropriate edge contributions; averaging
a gradient before squaring is generally different from averaging its square.
These choices are explicit in
\href{../src/turbulence/base/eddy_viscosity.cpp}{the shared SGS operators}.

For example, two neighboring shear samples $+a$ and $-a$ have a zero mean but
a mean square $a^2$. A closure that first averages them and then squares would
diagnose no deformation contribution at all. This is why matching array sizes
does not justify changing the staggered averaging order. The physical
interpretation is also limited: eddy viscosity parameterizes unresolved
transport at a filter scale; it is not molecular viscosity and does not
resolve the omitted eddies. Refining the grid while retaining a grid-derived
$\Delta$ changes both the resolved field and its unresolved closure.

Stable stratification can suppress turbulence. With the declared buoyancy
frequency squared $N^2$, the versioned Lilly-equilibrium branch uses a strain
factor equivalent to
\[
 |S|_{\mathrm{eff}}=
 \sqrt{\max(0,|S|^2-N^2/\mathrm{Pr}_t)}.
\]
Positive $N^2$ reduces mixing; negative $N^2$ increases the available scale.
Nimbus evaluates the expression with separate stable/unstable numerical forms
and retains a distinct older Richardson-correction option. Those options must
not be called identical. Moist versus dry buoyancy, coefficient caps and the
selected directional mode are part of the closure contract. The static
vertical-only tracked scope remains Engineering incomplete; the presence of
these formulas does not authorize every PBL, terrain or accelerator combination.

\section{Dynamic Smagorinsky diagnoses a bounded coefficient}

A dynamic model uses information at a second filter scale rather than prescribing
one coefficient everywhere. Nimbus reconstructs scalar-center fields, applies
a density-weighted Favre test filter and forms a resolved stress difference
$L_{ij}$ and a model difference $M_{ij}$. With a local least-squares assumption,
the coefficient multiplying the grid-scale model is estimated from their
contraction. The retained implementation averages horizontally at each height:
\[
 C(k)=\operatorname{clip}\!\left(
 \frac{\sum_{i,j}L_{ab}M_{ab}}{\sum_{i,j}M_{ab}M_{ab}},
 0,C_{s,\max}^2\right),\qquad C_s(k)=\sqrt{C(k)}.
\]
Repeated tensor indices denote the symmetric-tensor contraction, including
off-diagonal contributions. If the denominator is too small, the coefficient
is zero. The nonnegative bound excludes SGS backscatter in this implementation;
it is a deliberate model restriction, not a general theorem about turbulence.

The test width is twice the grid filter width. Favre filtering uses
$\widehat{\rho_d\phi}/\widehat{\rho_d}$, so ordinary averaging is not an
equivalent replacement in a variable-density field. Horizontal periodic
wrapping and bounded vertical sampling are also part of the retained filter.
The current owner requires its flat, periodic, C-grid, 3-D CPU domain and
rejects a required-native request. It is not a generic dynamic model for all
boundaries or mapped geometries. See
\href{../src/turbulence/schemes/smagorinsky/dynamic_smagorinsky.cpp}{dynamic
coefficient construction} and
\href{../tests/physics/test_turbulence_sgs.cpp}{SGS reference tests}.

Horizontal averaging is a substantive assumption. It supplies one coefficient
for a height plane and suppresses fluctuations in the least-squares ratio,
but can blend a strongly sheared storm core with a quieter environment. A
pointwise or Lagrangian average would define a different model, not simply a
faster reduction. The current test filter has separable three-point weights
$(1/4,1/2,1/4)$; its finite vertical edge treatment is distinct from a periodic
three-dimensional filter. The dynamic branch also rejects a separately
selected stability correction rather than composing two coefficient models
without a declared derivation.

\section{Prognostic TKE supplies a second modeled energy reservoir}

Let $e$ denote unresolved turbulent kinetic energy per unit dry mass, with units
m$^2$ s$^{-2}$. Its closure has the schematic material form
\[
 \frac{De}{Dt}=P_s+P_b-\epsilon
       +\frac{1}{\rho_d}\nabla\!\cdot(\rho_dK_e\nabla e).
\]
Shear production transfers energy through deformation; buoyancy production can
create or suppress turbulence; dissipation removes modeled TKE. Resolved
advection and source application still belong to the coupled timestep. A TKE
equation alone does not establish a conservative exchange with resolved kinetic
or internal energy.

Nimbus's Deardorff--Moeng-style branch uses $\nu_t=0.1\ell\sqrt e$. With its
stability correction active, the stable mixing length is limited by
$\ell=\min(\Delta,0.76\sqrt e/N)$ for $N>0$; neutral/unstable conditions use
the grid length, and disabling the correction retains that length directly.
For positive $e$, its retained dissipation is
\[
 \epsilon=\left(0.19+0.51\frac{\ell}{\Delta}\right)
             \frac{e^{3/2}}{\ell}.
\]
The zero-energy case has its own zero-source treatment. This avoids interpreting
the displayed quotient as permission to divide by a vanished length.

Scalar transport also depends on the length ratio. Before configured caps,
\[
\begin{aligned}
 K_\theta&=\nu_t\left[1+
        (\mathrm{Pr}_t^{-1}-1)\ell/\Delta\right],\\
 K_q&=\nu_t\left[1+
        (\mathrm{Sc}_t^{-1}-1)\ell/\Delta\right],\qquad K_e=2\nu_t.
\end{aligned}
\]
Thus replacing all three by one diffusivity changes the closure. Stress work,
TKE production and scalar diffusion use physical geometry and the admitted
discrete locations. The
\href{../src/turbulence/schemes/tke/tke.cpp}{TKE source} and
\href{../tests/physics/test_tke_scheme.cpp}{tests} define the actual branch,
including finite input, positivity and native-path requirements.

The stable length has the expected dimensions: $\sqrt e$ is a velocity and
$N$ an inverse time, so their ratio is a length. As stability reduces
$\ell/\Delta$, the scalar-to-momentum ratios above approach one. This is
different from applying a constant Prandtl number after computing viscosity.
The dissipation coefficient also varies with the length ratio, so changing
only the viscosity length would leave an inconsistent closure. Whether the
removed TKE is returned to another represented energy reservoir must be checked
in the actual coupled source/energy owner; the sign of $\epsilon$ alone does
not prove that exchange is implemented conservatively.

\section{PBL columns and surface exchange operate at another modeling scale}

The PBL choices are none, slab, YSU and MYNN. None supplies no PBL forcing; it
does not disable separately selected resolved-grid SGS. Active PBL models
represent vertical column transport and surface exchange with different
assumptions from a three-dimensional LES stress. Combining them requires an
explicit division of transport to avoid counting the same unresolved exchange
twice.

A useful column notation makes the nonlocal assumption visible:
\[
 J_\phi=-\rho_dK_\phi(\partial_z\phi-\gamma_\phi)
              +J_{\phi,\mathrm{entrainment}}.
\]
$\gamma_\phi$ has the units of a vertical gradient. For theta, $J_\theta$
has units kg K m$^{-2}$ s$^{-1}$ and still needs the declared heat conversion
to become W m$^{-2}$. For a water mixing ratio, the analogous flux is a water
mass flux. A positive countergradient contribution can transport a scalar
against its local gradient because it represents nonlocal eddies. Consequently,
one cannot infer the complete flux sign from $K>0$ alone. This schematic form
organizes the coupling; each PBL family supplies its own terms and limits.

The slab owner retains a mixed-layer state and diagnoses a bounded depth $h$.
Its column heights are measured above local ground. Weighting respects partial
layer overlap:
\[
 m_k=\rho_{d,k}\max\!\left(0,
       \min(h,z_{k+1/2})-z_{k-1/2}\right),\qquad
 \overline\phi_m=\frac{\sum_km_k\phi_k}{\sum_km_k}.
\]
These are masses per horizontal area. Mixing and surface forcing are applied
using complete cell masses so a partly occupied top layer does not receive a
full layer's forcing. Stored mixed-layer quantities and realized fluxes must
remain consistent with the resolved column. The word slab here does not imply
a prognostic soil reservoir. See
\href{../src/boundary_layer/schemes/slab/slab.cpp}{slab mixing}.

YSU supplies nonlocal transport and explicit entrainment, with physical
background in
\href{https://journals.ametsoc.org/view/journals/mwre/134/9/mwr3199.1.xml}{Hong,
Noh and Dudhia (2006)}. Nimbus's interior diffusivity profile has the
characteristic $K\sim\kappa w_m z(1-z/h)^2$ structure before its stability and
bounding choices. A scalar flux contains a downgradient part and, where
selected, a countergradient correction; entrainment adds a distinct top
exchange. The implemented temperature, moisture and momentum terms are not
all identical: for example, its moisture flux does not simply reuse the
temperature countergradient term. The
\href{../src/boundary_layer/schemes/ysu/ysu.cpp}{retained YSU owner} determines
the exact coefficients and interval solve.

MYNN uses prognostic TKE and stability functions, with background in
\href{https://www.jstage.jst.go.jp/article/jmsj/87/5/87_5_895/_article/-char/en}{Nakanishi
and Niino (2009)}. In Nimbus's level-2.5 notation, $q^2=2e$ is a velocity-variance
quantity, not water mixing ratio. The diffusivities have the form
\[
 K_m=\ell\sqrt{2e}\,S_m,
 \qquad K_h=\ell\sqrt{2e}\,S_h,
\]
with stability functions depending on shear, buoyancy and the mixing-length/TKE
state. The retained Helfand--Labraga correction limits unrealizable states;
invalid denominators or coefficients reject rather than silently producing
an arbitrary mixing rate. MYNN's column TKE is a closure state with its own
solve, not permission to add an independent LES TKE source twice. See
\href{../src/boundary_layer/schemes/mynn/mynn.cpp}{MYNN} and
\href{../src/boundary_layer/schemes/mynn/mynn_stability.cpp}{stability functions}.

\section{Surface-flux signs, units and prescribed skin state}

Let $\boldsymbol U$ be near-surface air velocity relative to the physical
surface, $U=|\boldsymbol U|$, and $\Pi=(p/p_{\rm ref})^{R_d/c_{pd}}$.
The retained bulk exchange takes the form
\[
\begin{aligned}
 \boldsymbol\tau&=-\rho_dC_dU\boldsymbol U,\\
 H&=\rho_dc_{pd}\Pi C_hU(\theta_s-\theta_a),\\
 E&=\rho_dC_eU(q_s-q_a).
\end{aligned}
\]
Stress opposes relative air motion; positive $H$ heats the atmosphere and
positive $E$ supplies vapor. Their units are N m$^{-2}$, W m$^{-2}$ and
kg m$^{-2}$ s$^{-1}$ respectively. $L_vE$ is a latent-energy flux, whereas $E$
itself is not. The Exner factor in $H$ converts a potential-temperature flux
to sensible-heat currency; omitting it changes the implemented exchange.

Neutral coefficients use separate momentum and heat roughnesses:
\[
 C_{d0}=\frac{\kappa^2}{\ln^2(z/z_{0m})},\qquad
 C_{h0}=\frac{\kappa^2}{\ln(z/z_{0m})\ln(z/z_{0h})}.
\]
The bulk branch applies retained Louis stability factors. The Monin--Obukhov
branch instead iterates stability-corrected logarithmic denominators using
$\zeta=z/L$ and the surface virtual-potential-temperature flux:
\[
 L=-\frac{u_*^3\theta_{v,a}}
          {\kappa g\,\overline{w'\theta_v'}_s},\qquad
 \overline{w'\theta_v'}_s\simeq
 \frac{H}{\rho_dc_{pd}\Pi}+0.61\theta_a\frac{E}{\rho_d}.
\]
An unstable positive virtual heat flux therefore gives negative $L$. Nimbus
retains explicit neutral limits, iteration limits and minimum-friction-velocity
handling; these are part of the algorithm, not an unbounded analytic law.

The minimum-$u_*$ treatment can increase momentum and heat/moisture transfer
coefficients together. It is therefore a modification of the exchange, not a
harmless display floor. Roughness must be below the reference sampling height,
and the neutral limit must remain finite. These checks are necessary even for
a flat surface: a logarithmic formula evaluated outside its geometric domain
does not become physical by clipping its final flux afterward.

Skin temperature and moisture are prescribed, or initialized from the initial
surface air when unset and subsequently held. There is no prognostic soil
heat/moisture or vegetation model behind them. On sloping ground, tangent wind,
traction and actual-versus-horizontal area must be transformed consistently.
See \href{../src/boundary_layer/base/surface_fluxes.cpp}{surface exchange}.
Neither radiative surface absorption nor a surface-flux receipt automatically
evolves a land temperature that the model does not own.

\section{Diffusion applies a closure; it does not define one}

The density-weighted scalar diffusion form is
$\rho_d\partial_t\phi=\nabla\!\cdot(\rho_dK\nabla\phi)$.
The coefficient can come from a prescribed numerical profile or a physical
closure. Explicit time integration has a diffusion stability restriction;
implicit vertical exchange solves a coupled column system over the interval.
Neither method makes an inappropriate coefficient physically justified.

Nimbus separates coefficient construction, transport/application and realized
source reports. An implicit solve's accepted flux need not equal an explicit
flux formed from the old gradient. PBL held sources retain interval-mean
forces and inventory rates so changing density does not accidentally reuse a
stale acceleration denominator. Momentum diffusion also changes resolved
kinetic energy, with a finite-step contribution beyond linear work. Appropriate
boundary work and TKE/thermal attribution must be assessed together.
The existing diffusion chapter derives the spatial and temporal operators;
the closure-specific question is which $K$, carrier, interval and coupled
energy owner feed them.

Some admitted paths defer transport-tendency construction until the selected
application owner consumes closure coefficients. That prevents an explicit
tendency from being applied and then implicitly applied again over the same
interval. It also makes coefficient lifetime significant: a matrix built from
one state is not automatically the matrix corresponding to a later gradient
or density. Changes between explicit, implicit and deferred paths need tests
of the accepted state and realized flux, including rejection, rather than a
comparison of coefficient arrays alone.

\section{Stochastic perturbations have covariance and conservation contracts}
\label{chaos-theory}

The chaos choices are none, initial-condition perturbations, boundary-layer
perturbations and full stochastic coupling. A simple temporal building block
is the retained AR(1) update
\[
 \xi_{n+1}=a\xi_n+\sqrt{1-a^2}\,\eta_n,
 \qquad a=\exp(-\Delta t/\tau).
\]
If independent innovations $\eta_n$ and the previous field each have zero mean
and unit variance, the new variance is $a^2+(1-a^2)=1$. That explains the
innovation coefficient. Finite-grid normalization, spatial filtering,
clipping and conditioning can change the realized distribution. The variance
identity does not establish the exact finite-domain distribution after those
additional transformations. The actual noise refresh interval must
be used; replacing it with an unrelated solver step changes decorrelation.

The spectral spatial filter multiplies Fourier components by
\[
 G(k_x,k_y)=\exp[-\tfrac12\{(L_xk_x)^2+(L_yk_y)^2\}].
\]
This attenuates short scales. The recursive alternative applies repeated
forward/backward smoothing with weights such as $\exp(-\Delta x/L_x)$.
It approximates a correlated field through a different discrete operation;
it is not identical to the spectral filter. Filter widths, boundary treatment
and variance normalization affect the resulting covariance. A configured
length is not automatically the measured autocorrelation length of every
clipped or relocated field. See
\href{../src/chaos/base/correlation_filter.cpp}{correlation filters} and
\href{../src/chaos/base/perturbation_field.cpp}{temporal updates}.

Seeds are combined with member, process/field and sequence identity. Restart
must preserve the retained noise generation and clock. Earth-correlated noise
uses its declared physical control-grid coordinates; a translated window must
not invent a new forcing realization solely because an Earth point acquires
a different array index. Different standard-library random-distribution
implementations are not automatically bitwise equivalent despite equal seeds.

A source multiplier of the form $1+\alpha\xi$, subject to the declared bounds,
changes a process. It must be applied to the associated transfers or conserved
column exchange consistently. Scaling water tendencies independently can break
their zero sum; scaling a sedimentation divergence without its lower-boundary
flux breaks precipitation accounting. Nimbus retains source/column ownership
and private-state checks for admitted stochastic couplings. The current
\href{../src/chaos/schemes/full_stochastic/full_stochastic.cpp}{full-stochastic
owner} is not a license to multiply every physical term by unrelated noise.
Ensemble spread is also not calibrated forecast uncertainty merely because
the perturbation variance is known.

Bounding a multiplier also changes its statistics: even if $\xi$ has zero mean,
the mean of a clipped $1+\alpha\xi$ need not be exactly one. Common process
identity and a conserved-transfer construction address a different question
from unbiased ensemble forcing; neither follows automatically from the AR(1)
variance calculation.

\section{Radar products are bulk forward operators}
\label{radar-diagnostics-theory}

Radar diagnostics translate modeled particles and wind into an observation-like
quantity. A Rayleigh liquid reference motivates the sixth-diameter moment
$Z=\int D^6n_v(D)\,dD$, with a volume spectrum and conversion to
mm$^6$ m$^{-3}$. Its logarithmic form is
\[
 Z_{\rm dBZ}=10\log_{10}\!\left(
       Z_e/(1\ \mathrm{mm^6\,m^{-3}})\right).
\]
Frozen-particle density, dielectric behavior and shape require further
assumptions. Summing component reflectivities must occur in linear space,
not by adding their dBZ values.

Nimbus's \texttt{fast\_da} and \texttt{psd\_moment} tiers retain concrete
bulk relationships. Without supplied number, a species path uses
$Z_{e,x}=A_x(\rho_dq_x)^{1.75}$. With supplied number, its moment proxy has
the form $Z_{e,x}=C_x(\rho_dq_x)^2/N_{v,x}$ with retained species factors.
The number used by that interface is interpreted per volume; a private number
per kg dry air needs conversion before being supplied. The tier can use the
mass-only relation when the corresponding number is absent. These are retained
proxies, not an arbitrary gamma-spectrum integral or a general electromagnetic
scattering solver. See
\href{../src/radar/schemes/reflectivity/reflectivity.cpp}{reflectivity tiers}.

Radial velocity projects reconstructed physical wind onto the vector from the
radar to the sample:
$V_r=\boldsymbol u\cdot\widehat{\boldsymbol r}$, with the configured bulk
scatterer fall correction when selected. C-grid face winds are reconstructed
at the declared scalar sample before projection. Terrain heights and Earth
origin determine the line of sight. The current \texttt{beam\_volume} option
uses a bounded local neighborhood average; it is not a complete antenna beam,
pulse-volume, attenuation or instrument-noise simulator. The
\href{../src/radar/schemes/velocity/velocity.cpp}{velocity owner} defines its
sampling and fall-speed bounds.

Differential reflectivity is $Z_{DR}=10\log_{10}(Z_H/Z_V)$ where the signals
are valid. Nimbus retains \texttt{simple}, \texttt{polarimetric\_fo} and
\texttt{zhang2021\_rain\_mp} paths. The bulk polarimetric path uses prescribed
shape/species relationships and publishes associated KDP and correlation
products within declared ranges; its name does not establish a T-matrix
calculation. The separate rain-MP path rejects frozen-species participation.
KDP is expressed in degrees per kilometre and the correlation coefficient is
dimensionless; these units do not make a bulk proxy a resolved beam-phase model.
Low/no-signal conventions and clipping are part of output interpretation.
The \href{../src/radar/schemes/zdr/zdr.cpp}{polarimetric source} and its tests
must accompany any claim about an admitted product set.

\section{Diagnostic integrals, storm objects and parcel paths}

Budget diagnostics integrate physical inventories and boundary exchanges,
using their versioned dry/mixture mass and phase-reference conventions. An
exact telescoping ledger records how a diagnostic changed; independent source,
work and boundary references are still needed to validate its physics.
Derived gradients and layer integrals require physical heights and finite
coverage, not a nominal index range.

The sounding diagnostic lifts a surface parcel dry-adiabatically to its LCL
and with a retained moist lapse-rate approximation above it. Vapor-corrected
virtual temperatures give buoyancy
$B=g(T_{v,p}-T_{v,e})/T_{v,e}$. Its CAPE is the positive buoyancy integral
between diagnosed LFC and EL; CIN is reported as a nonnegative magnitude of
the negative integral below LFC. Nimbus uses discrete layer integration and
finite column coverage. With no LFC it reports unavailable crossing heights
and a no-LFC interpretation with zero CIN, not evidence that inhibition is
physically absent. A top-limited EL likewise needs its finite-column context.
See \href{../src/init/sounding/diagnostics.cpp}{sounding diagnostics}.

Storm detection uses resolved updraft, condensate and rotation evidence to form
feature candidates. Temporal association, branching and ambiguity rules select
a track rather than identifying a uniquely defined physical object from one
maximum. Updraft helicity has the dimensional form $\int w\zeta\,dz$ over its
declared updraft/height support; storm-relative helicity uses the signed
storm-relative hodograph convention. Their values depend on physical layer
coverage and the selected storm motion. The
\href{../src/diagnostics/storm_tracker.cpp}{tracker} retains history because
continuity cannot be reconstructed from a single frame.

Air-parcel trajectories solve $d\boldsymbol x/dt=\boldsymbol u(\boldsymbol x,t)$
under the model's interpolation and frame assumptions. Nimbus prepares a
predicted midpoint from the beginning wind, samples there, and combines that
midpoint sample with the end-state midpoint wind during completion. This
differs from a streamline in one frozen velocity field. Stagger-aware sampling,
physical height, vertical clamping, horizontal departure and Earth-frame output
are explicit choices; the paths are not inertial precipitation trajectories.
\href{../src/diagnostics/trajectories/trajectories.cpp}{Trajectory continuation}
retains parcels and sample cadence through restart and movement.

These diagnostics test different aspects of a simulation. Matching rainfall,
CAPE or a storm track does not uniquely validate all closures that produced
them. Bounded formula, coupling and failure checks establish their stated
implementation scope; long campaigns and independent scientific review remain
separate from those checks.
